% =====================================================================
%  CARNET D'ENTRAÎNEMENT — CHIMIE — FICHE 3 — THÈME 2
%  Composition : atomes et ions — NIVEAU FACILE — ÉNONCÉS
% =====================================================================
\documentclass[11pt,a4paper]{article}

\usepackage[utf8]{inputenc}
\usepackage[T1]{fontenc}
\usepackage{lmodern}
\usepackage[french]{babel}

\usepackage{amsmath,amssymb,mathtools}
\usepackage{xcolor}
\usepackage{colortbl}
\usepackage{tikz}
\usetikzlibrary{arrows.meta,positioning,calc}
\usepackage[most]{tcolorbox}
\usepackage{enumitem}
\usepackage{booktabs}
\usepackage{array}
\usepackage[a4paper,margin=2.1cm,headheight=26pt,headsep=9pt]{geometry}
\usepackage{fancyhdr}
\usepackage[hidelinks]{hyperref}

\definecolor{primary}{RGB}{124,78,140}
\definecolor{primarydark}{RGB}{74,44,92}
\definecolor{defcol}{RGB}{0,114,178}
\definecolor{thmcol}{RGB}{0,140,120}
\definecolor{attncol}{RGB}{200,40,55}

\newcommand{\nuc}[3]{\prescript{#1}{#2}{\mathrm{#3}}}
\newcommand{\pp}{p^{+}}
\newcommand{\ee}{e^{-}}
\newcommand{\nz}{n^{0}}

\pagestyle{fancy}
\fancyhf{}
\fancyhead[L]{\small\textcolor{primarydark}{\textbf{Fiche 3 — Thème 2}} : Composition des atomes et des ions}
\fancyhead[R]{\small\textcolor{primarydark}{\textsc{Facile — Énoncés}}}
\fancyfoot[C]{\small\thepage}
\renewcommand{\headrulewidth}{0.8pt}
\renewcommand{\headrule}{\hbox to\headwidth{\color{primary}\leaders\hrule height \headrulewidth\hfill}}

\newcounter{exo}
\newcommand{\exo}[1][]{\stepcounter{exo}\par\medskip%
  \noindent{\color{primary}\rule[-2pt]{3.5pt}{15pt}}\ \textbf{\color{primarydark}\large Exercice \theexo}%
  \ifx\relax#1\relax\relax\else\ \textmd{\normalsize\color{black!65}— \itshape #1}\fi\par\nopagebreak\smallskip}

\setlist{leftmargin=1.7em,topsep=2pt,itemsep=3pt}

\begin{document}

\begin{tikzpicture}
\node[rectangle,rounded corners=6pt,fill=primary,text=white,
      minimum width=\textwidth,minimum height=1.35cm,align=center,inner sep=6pt]
  {\Large\bfseries Thème 2 — Composition des atomes et des ions\\[2pt]
   \normalsize\mdseries Niveau Facile — Énoncés};
\end{tikzpicture}

\vspace{3pt}
{\small\itshape Rappels : $Z = $ nombre de protons ; $A = $ nombre de nucléons ; nombre de neutrons
$= A - Z$ ; atome neutre $\Rightarrow$ nombre d'électrons $= Z$.}

% ---------------------------------------------------------------
\exo[lire une notation symbolique]
On considère l'atome de chlore $\nuc{35}{17}{Cl}$. Donne :
\begin{enumerate}[label=\alph*)]
  \item son numéro atomique $Z$ et son nombre de masse $A$ ;
  \item son nombre de protons et son nombre de neutrons ;
  \item son nombre d'électrons, sachant qu'il est neutre.
\end{enumerate}

% ---------------------------------------------------------------
\exo[compter dans un atome neutre]
L'atome de calcium a pour numéro atomique $Z = 20$ et pour nombre de masse $A = 40$. Combien
possède-t-il de protons, de neutrons et d'électrons ?

% ---------------------------------------------------------------
\exo[écrire la notation]
Un atome de sodium a $Z = 11$ et $A = 23$, symbole $\mathrm{Na}$. Écris sa notation symbolique
$\nuc{A}{Z}{X}$ en plaçant correctement $Z$ et $A$. Lequel des deux se met \emph{en haut} ? lequel
\emph{en bas} ?

% ---------------------------------------------------------------
\exo[que représentent $Z$ et $A$ ?]
Dans la notation $\nuc{A}{Z}{X}$ :
\begin{enumerate}[label=\alph*)]
  \item que représente le nombre $Z$ ?
  \item que représente le nombre $A$ ?
  \item comment obtient-on le nombre de neutrons à partir de $Z$ et $A$ ?
\end{enumerate}

% ---------------------------------------------------------------
\exo[remonter à $Z$ et $A$]
Un atome possède $8$ protons et $8$ neutrons dans son noyau.
\begin{enumerate}[label=\alph*)]
  \item Quel est son numéro atomique $Z$ ?
  \item Quel est son nombre de masse $A$ ?
  \item Combien a-t-il d'électrons s'il est neutre ?
\end{enumerate}

\end{document}
